\chapter{Introducción}
\label{ch:introduccion}

Este informe técnico documenta el desarrollo y la resolución del Ejercicio Práctico N° 1 del Trabajo Práctico N° 2
de la materia Técnicas Digitales II. Las actividades propuestas se llevan a cabo sobre la placa de desarrollo
BluePill, la cual incorpora un microcontrolador STM32F103C8T6 basado en el núcleo ARM Cortex-M3. El enfoque de
programación implementado es Bare Metal puro, interactuando de forma directa con los registros de los periféricos a
través de punteros de memoria.

La guía de trabajos prácticos plantea una serie de cinco ejercicios orientados a comprender el funcionamiento y
configuración de diversos módulos del microcontrolador:

\begin{enumerate}
    \item Ejercicio Práctico N° 1 (Interrupciones Externas y el NVIC): Configuración del controlador de
    interrupciones vectorizadas y anidadas (NVIC) para gestionar eventos asíncronos provenientes de un pin de
    entrada (botón) mediante interrupción externa (EXTI), modificando el estado de un LED.
    \item Ejercicio Práctico N° 2 (Conversión Analógica-Digital): Configuración del ADC para la lectura
    del sensor de temperatura interno y de un canal analógico externo.
    \item Ejercicio Práctico N° 3 (Migración a CMSIS y Gestión de Tiempos): Incorporación del estándar
    CMSIS para mejorar la legibilidad y uso del temporizador del sistema (SysTick) para generar retardos precisos.
    \item Ejercicio Práctico N° 4 (Comunicación Serie UART): Establecimiento de comunicación
    bidireccional con una PC mediante el periférico USART para la transmisión de estados y recepción de comandos.
    \item Ejercicio Práctico N° 5 (Control de Potencia con Timers): Generación de señales de modulación por
    ancho de pulso (PWM) empleando los temporizadores de propósito general para variar la intensidad lumínica.
\end{enumerate}

Cabe destacar que el presente documento aborda de manera exclusiva la resolución del Ejercicio Práctico N° 1. Los
ejercicios subsiguientes (del 2 al 5) se desarrollarán de forma progresiva en futuros informes. El objetivo principal
de esta primera etapa consiste en comprender la tabla de vectores de interrupción, el funcionamiento de la línea EXTI
y la habilitación de interrupciones en el NVIC sin utilizar bibliotecas de alto nivel.
